\section{All integer kernels: a terminating harmonic recurrence}
\label{bmh:sec:kernels}

\subsection{Laurent transport on the compact interval}
Put
\[
 y(t)=\log\frac{t}{1-t},\qquad
 L_j(t)=\Li_j\!\left(-\frac{t}{1-t}\right).
\]
The identities required by the recurrence are
\begin{equation}\label{bmh:eq:Lderivative}
 L_0(t)=-t,\qquad
 L_j'(t)=\frac{L_{j-1}(t)}{t(1-t)}\quad(j\ge1),\qquad
 y'(t)=\frac1{t(1-t)}.
\end{equation}
In particular the zero-order value is $-t$, not $-t/(1-t)$.

For a Laurent polynomial $R(t)=\sum_jr_jt^j$, define
\begin{equation}\label{bmh:eq:SR}
 c_R=r_{-1},\qquad
 S_R(t)=\sum_{j\ne-1}r_j\frac{t^{j+1}-1}{j+1},\qquad
 K_R(t)=\frac{S_R(t)}{t(1-t)}.
\end{equation}
Since $S_R(1)=0$, $K_R$ is again a Laurent polynomial. Let
\[
 F_R(m,n;h)=\int_0^1 R(t)y(t)^h L_m(t)L_n(t)\dd t,
\]
whenever the displayed integral is ordinary and convergent.

\begin{theorem}[Finite integer-kernel reduction]\label{bmh:thm:kernel}
For $m,n\ge1$ and every integrable instance reached from
$R=R_{N,a}$ with $-1\le a\le N-1$ integer,
\begin{align}
 F_R(m,n;h)={}&c_R I_{m,n}^{(h)}
 -F_{K_R}(m-1,n;h)-F_{K_R}(m,n-1;h)\notag\\
 &-hF_{K_R}(m,n;h-1),\label{bmh:eq:kernel-rec}
\end{align}
where the last term is omitted when $h=0$. Whenever an index reaches
zero, absorb $L_0=-t$ into the rational kernel and remove that factor.
The one-factor recurrence is
\begin{equation}\label{bmh:eq:one-rec}
 F_R(n;h)=c_R J_n^{(h)}-F_{K_R}(n-1;h)-hF_{K_R}(n;h-1),
\end{equation}
with
\begin{equation}\label{bmh:eq:one-base}
 J_n^{(h)}=-h!\sum_{j=0}^{\lfloor h/2\rfloor}
 e_{2j}\frac{(n)_{h+1-2j}}{(h+1-2j)!}\,
            \zeta(n+h+1-2j).
\end{equation}
The final zero-factor bases are the Gamma--harmonic moments in
Theorem~\ref{bmh:thm:beta-base}. These rules terminate and evaluate every
$\partial_a^h M_N(a;m,n)$ at every admissible integer $a$.
\end{theorem}
\begin{proof}
Write $R=c_R/t+S_R'$ and integrate the second term by parts. Use
\eqref{bmh:eq:Lderivative} on each factor and on $y^h$ to get
\eqref{bmh:eq:kernel-rec}. The boundary at one vanishes because
$S_R(1)=0$ and each $L_j$ and $y$ is at most logarithmic there. At zero,
the initial bilinear kernel has pole order at most two. Two positive-order
factors are $O(t^2)$, and $S_R$ has pole order at most one, so the boundary
vanishes even with logarithmic powers. The recurrence preserves this
property. Absorbing a zero-order factor multiplies the kernel by $-t$;
with one factor left the pole order is at most one, and with no factors
left the kernel is an ordinary polynomial. Thus all reached boundary
terms and integrals have the asserted ordinary meaning.

The same argument proves~\eqref{bmh:eq:one-rec}. Its base
$J_n^{(h)}=\int_0^1y^hL_n(t)\dd t/t$ is the one-factor resonance from
\cite{MLD}. To check its normalization directly, use the already proved
one-factor transform
\[
 J_1(a,n;1)=\pi\csc(\pi a)\{\zeta(n)-\zeta(n;1-a)\}
\]
and its cancelled Taylor series at $a=0$. At $n=1$, replace the zeta
difference by $\psi(1-a)-\psi(1)$ before expansion. This gives exactly
\eqref{bmh:eq:one-base}, including its minus sign.

Every recursive call reduces the sum of the positive orders and $h$ by
one, or removes a zero-order factor. Polynomial decomposition is finite.
The process therefore terminates without solving an infinite recurrence.
Finally~\eqref{bmh:eq:compact-kernel} identifies the initial integral
with the desired Mellin moment.
\end{proof}

\begin{corollary}[Representation class at integer parameters]
At $h=0$, every unit-scale $M_N(a;m,n)$ with $a=-1,0,\ldots,N-1$ is a
rational linear combination of rational numbers, ordinary zeta values,
and convergent double zeta values, of weights at most $m+n+1$.
For $a\ge1$ the weight bound improves to $m+n$. With $h$ logarithms the
bounds become $m+n+h+1$ and $m+n+h$, respectively. The answers belong
to the algebra generated by single and double zeta values; after stuffle,
an additive depth bound of three is sufficient for the formulas produced
here.
\end{corollary}
\begin{proof}
Use Theorem~\ref{bmh:thm:jets}, the one-factor base, and the elementary
bases below in the terminating recurrence. The sharper weight bounds
are also immediate from~\eqref{bmh:eq:top-weight}. In the coefficient
formulas the only extra products are even zeta factors multiplying a
single or double value, or the two-zeta corner terms. Products of even
zeta values are rational multiples of a single even zeta value. Thus
stuffle requires at most three nested indices.
\end{proof}

\subsection{The Gamma base and generalized harmonic numbers}
\begin{theorem}[Beta moments and Bell polynomials]\label{bmh:thm:beta-base}
For integers $d,h\ge0$,
\begin{align}
 B_{d,h}&:=\int_0^1t^d\log^h\frac{t}{1-t}\dd t\notag\\
 &=\left.\frac{\dd^h}{\dd v^h}B(d+1+v,1-v)\right|_{v=0}\notag\\
 &=\frac{h!}{d+1}[v^h]\left\{
       \frac{\pi v}{\sin\pi v}\prod_{j=1}^d\left(1+\frac vj\right)
       \right\}.\label{bmh:eq:beta-base}
\end{align}
Equivalently, $B_{d,h}=Y_h(\kappa_1,\ldots,\kappa_h)/(d+1)$, where
$Y_h$ is the complete exponential Bell polynomial and
\begin{align}
 \kappa_1&=H_d,\label{bmh:eq:kappa1}\\
 \kappa_r&=(r-1)!\left\{
 (1+(-1)^r)\zeta(r)+(-1)^{r+1}H_d^{(r)}\right\},\qquad r\ge2.
 \label{bmh:eq:kappas}
\end{align}
\end{theorem}
\begin{proof}
Differentiate the ordinary beta integral near $v=0$; an integrable
power-logarithmic majorant justifies the derivatives. The Gamma quotient
and reflection identity give
\[
 B(d+1+v,1-v)
 =\frac1{d+1}\Gamma(1+v)\Gamma(1-v)
           \prod_{j=1}^d(1+v/j)
 =\frac1{d+1}\frac{\pi v}{\sin\pi v}\prod_{j=1}^d(1+v/j).
\]
This proves~\eqref{bmh:eq:beta-base}. Taking logarithmic derivatives
of the product gives~\eqref{bmh:eq:kappa1}--\eqref{bmh:eq:kappas}.
The relation between derivatives of an exponential and the complete
Bell polynomials proves the last assertion. These beta and Gamma tools
are classical; their role here is to close the bilinear recurrence.
\end{proof}

\subsection{Examples with different denominator orders}
The recurrence gives, without numerical fitting,
\begin{align}
 M_2(1;1,1)&=2,\label{bmh:eq:kernel-example1}\\
 M_3(1;2,1)&=-\frac54+\frac14\zeta(2)+\zeta(3),\label{bmh:eq:kernel-example2}\\
 M_3(2;2,2)&=\zeta(2)+2\zeta(3)+\frac{17}{4}\zeta(4),\label{bmh:eq:kernel-example3}\\
 M_2(0;2,1)&=\frac{17}{4}\zeta(4)-\zeta(2)-2\zeta(3),\label{bmh:eq:kernel-example4}\\
 \left.\partial_a M_2(a;2,2)\right|_{a=1}
 &=15\zeta(5)+12\zeta(2)\zeta(3).\label{bmh:eq:kernel-example5}
\end{align}
Their exact expressions are regenerated by the supplied symbolic code.
The two-factor weight bound does not imply homogeneity when the rational
kernel has polynomial parts: the lower-weight terms in these examples
are expected, not anomalous.
